% =====================================================================
% บทที่ 8 การแก้ไขปัญหาเบื้องต้น (ผู้อ่าน: ทุกกลุ่ม)
% รวบรวมจากปัญหาที่พบจริงระหว่างทดสอบระบบกับผู้ใช้ (UAT) และระหว่างติดตั้ง
% =====================================================================
\label{ch:trouble}

\begin{guide}
รวบรวมจากปัญหาที่พบจริงระหว่างทดสอบระบบกับผู้ใช้ (UAT) และระหว่างติดตั้งบนเครื่องใหม่
\begin{itemize}
  \item ตารางอาการ--สาเหตุ--วิธีแก้ แยกปัญหาของผู้ใช้ออกจากปัญหาการติดตั้งและดูแลระบบ
  \item คัดลอกข้อความแจ้งเตือนให้ตรงกับที่ระบบแสดงจริง ผู้อ่านจะค้นหาข้อความนั้นในคู่มือได้
  \item คำถามที่พบบ่อย และช่องทางติดต่อขอความช่วยเหลือ พร้อมข้อมูลที่ควรแจ้งมา
\end{itemize}
\end{guide}

\section{ปัญหาที่ผู้ใช้พบบ่อย}

\begin{troubleshooting}[ปัญหาที่ผู้ใช้พบบ่อยและวิธีแก้ไข][tab:troubleshoot]
\problem{สมัครสมาชิกไม่ได้ ระบบแจ้งว่าอีเมลไม่ถูกต้อง}
        {ใช้อีเมลที่ไม่ได้ลงท้ายด้วย \cmd{@kku.ac.th}}
        {สมัครด้วยอีเมลของมหาวิทยาลัย}
\problem{Compile แล้วไม่ได้ไฟล์ PDF}
        {โค้ด \LaTeX{} ผิดรูปแบบ เช่น ปิด \cmd{\begin{...}} ไม่ครบ หรือวงเล็บปีกกาไม่ครบคู่}
        {อ่านบรรทัดที่ผิดใน COMPILE LOG แล้วแก้ไข}
\problem{ขึ้นข้อความ ``กำลังแปลผลอยู่''}
        {สมาชิกคนอื่นในโครงงานกำลัง Compile}
        {รอสักครู่แล้ว Compile ใหม่}
\problem{จุดสถานะที่แถบเครื่องมือเป็นสีแดง}
        {WebSocket ไม่ได้เชื่อมต่อ ระบบทำงานแบบออฟไลน์}
        {รีเฟรชหน้าเว็บ และบันทึกงานด้วย \keys{Ctrl+S}}
\problem{เห็นตัวอักษรเพิ่มขึ้นใน Editor เอง}
        {สมาชิกคนอื่นกำลังแก้ไขไฟล์เดียวกัน}
        {ดูรายชื่อผู้ที่ออนไลน์ที่แถบเครื่องมือ ระบบรวมการแก้ไขให้อัตโนมัติ}
\problem{ภาพที่แทรกไม่แสดงใน PDF}
        {ยังไม่ได้อัปโหลดไฟล์ภาพ หรือชื่อไฟล์ใน \cmd{\includegraphics} ไม่ตรง}
        {อัปโหลดภาพที่แฟ้มข้อมูล แล้วแก้ชื่อไฟล์ให้ตรงกัน}
\end{troubleshooting}

\section{ปัญหาการติดตั้งและดูแลระบบ}

\begin{troubleshooting}[ปัญหาการติดตั้งและดูแลระบบ][tab:troubleshoot-admin]
\problem{\cmd{EADDRINUSE: port 3001}}
        {มีโปรแกรมอื่นใช้พอร์ต 3001 อยู่}
        {ปิดโปรแกรมเดิม หรือเปลี่ยน \cmd{PORT} ในไฟล์ \filename{server/.env}}
\problem{\cmd{xelatex: command not found} หรือ Compile ไม่ได้ PDF ทุกโครงงาน}
        {ยังไม่ได้สร้างอิมเมจ TeX Live หรือไม่ได้ติดตั้ง TeX Live}
        {สร้างอิมเมจตามขั้นตอนที่ 3 ในหัวข้อ~\ref{sec:install-steps}}
\problem{\cmd{npm install} ล้มเหลวที่แพ็กเกจ SQLite}
        {เครื่องไม่มีเครื่องมือ build ของ native module}
        {Linux: \cmd{sudo apt install build-essential python3} \newline Windows: ติดตั้ง Visual Studio Build Tools}
\problem{เบราว์เซอร์แจ้ง \cmd{CORS error}}
        {\cmd{CLIENT_URL} ไม่ตรงกับที่อยู่ของส่วนหน้า}
        {แก้ \cmd{CLIENT_URL} ใน \filename{server/.env} แล้วเริ่มส่วนหลังใหม่}
\problem{ผู้ใช้แก้ไขพร้อมกันไม่ได้หลังติดตั้งบนเครื่องแม่ข่าย}
        {Nginx ไม่ได้ส่งต่อ WebSocket ที่ \cmd{/collab}}
        {เพิ่ม \cmd{proxy_set_header Upgrade} และ \cmd{Connection} ตามไฟล์ในหัวข้อ~\ref{sec:install-steps}}
\problem{\cmd{database is locked} (ระหว่างพัฒนา)}
        {มีหลายโปรแกรมเปิดไฟล์ SQLite พร้อมกัน}
        {ปิดส่วนหลังที่เปิดค้างทั้งหมด แล้วเริ่มใหม่}
\end{troubleshooting}

\section{คำถามที่พบบ่อย}

\begin{faq}
\question{ทำไมข้อเสนอโครงงานถูกส่งกลับ}
อาจารย์ที่ปรึกษาส่งกลับเพื่อให้แก้ไข ให้อ่านความเห็นของอาจารย์ที่หน้าข้อเสนอ แก้ไขตามความเห็น แล้วส่งใหม่

\question{เพิ่มรูปภาพในเอกสารอย่างไร}
อัปโหลดรูปที่แฟ้มข้อมูล แล้วใช้ปุ่ม \menu{แทรก > ภาพ} เพื่อแทรกคำสั่ง \cmd{\includegraphics}
จากนั้นแก้ชื่อไฟล์ในคำสั่งให้ตรงกับไฟล์ที่อัปโหลด

\question{อ้างอิงเอกสารอย่างไร}
เพิ่มรายการในไฟล์ \filename{references.bib} แล้วใช้คำสั่ง \cmd{\cite{key}} ในเนื้อหา

\question{ลืมรหัสผ่านต้องทำอย่างไร}
ติดต่อผู้ดูแลระบบตามช่องทางในหัวข้อ~\ref{sec:contact} ผู้ดูแลระบบจะตั้งรหัสผ่านชั่วคราวให้
(หัวข้อ~\ref{sec:users}) แล้วเปลี่ยนรหัสผ่านใหม่ที่หน้าโปรไฟล์หลังเข้าสู่ระบบ
\end{faq}

\section{การติดต่อขอความช่วยเหลือ}
\label{sec:contact}

หากพบปัญหาที่ไม่มีในคู่มือ ติดต่อผู้ดูแลระบบหรือผู้จัดทำได้ที่
สาขาวิชาวิทยาการคอมพิวเตอร์ วิทยาลัยการคอมพิวเตอร์ มหาวิทยาลัยขอนแก่น
โดยแจ้งชื่อผู้ใช้ เวลาที่พบปัญหา และข้อความแจ้งเตือนหรือภาพหน้าจอประกอบ
